\subsection{线性表示}

定义 1. —— 设 M 是一个 K-模。G 在 M 中的线性表示是指从 G 到线性群 GL(M) 的一个群同态（II, §1, No.~2, 第 195 页）。

设 \(\pi\) 是 G 在 K-模 M 中的一个线性表示。我们也说偶 \((M,\pi)\) 是 G 的一个线性表示。当 K 非零且 M 是自由 K-模时，M 的维数称为表示 \(\pi\) 的次数（或维数）。

回顾一下，典范映射 \(g \mapsto e_{g}\)（从 G 到群 G 的代数 K{[}G{]} 的典范映射）是 K-模 K{[}G{]} 的一组基。滥用记号，我们也把元素 \(e_{g}\)（K{[}G{]} 的元素）记作 g。给定 G 的一个线性表示 \(\pi : G \to \mathbf{GL}(M)\)，存在唯一的从 K{[}G{]} 到 \(\operatorname{End}_{\mathrm{K}}(\mathrm{M})\) 的 K-代数同态延拓 \(\pi\)（III, §2, No.~6, 第 447 页，例）；我们也把这个延拓记作 \(\pi\)。它是代数 K{[}G{]} 在 M 中的一个表示（VIII, 第 373 页，定义 1），从而 M 是一个 K{[}G{]}-模。反之，每个线性表示 \(\rho\)（代数 K{[}G{]} 在 M 中的线性表示）都定义 G 在 M 中的一个线性表示，它由 \(g \mapsto \rho(g)\) 给出。

我们将不受拘束地把有关 K{[}G{]}-模结构的术语用于群 G 的线性表示。例如，我们说到子表示、单表示、表示的直和，等等。设 \((\mathrm{M},\pi)\) 与 \((\mathrm{M}^{\prime},\pi^{\prime})\) 是 G 的线性表示；用 \(\operatorname{Hom}_{\mathrm{G}}(\pi,\pi^{\prime})\) 表示 K-模 \(\operatorname{Hom}_{\mathrm{K[G]}}(\mathrm{M},\mathrm{M}^{\prime})\)（从 M 到 M' 的 K{[}G{]}-模同态所成的 K-模）。\(^{(1)}\)

映射 \(f\)（从 \(G\) 到 \(K\) 的映射）称为中心函数，如果 \(f(gg') = f(g'g)\) 对每一对 \((g, g')\)（\(G\) 中元素的每一对）成立；等价地，\(f(ghg^{-1}) = f(h)\) 对每一对元素 \(g, h\)（\(G\) 的元素对）成立。因此，中心函数就是 \(G\) 上在每个共轭类上取常值的函数。它们构成从 \(G\) 到 \(K\) 的映射空间的一个子模，记作 \(\mathcal{Z}_K(G)\)。当 \(G\) 有限时，K-代数 \(\mathcal{Z}_K(G)\) 是一个自由 K-模，其维数为 \(G\) 的共轭类的个数。中心 \(Z(K[G])\)（代数 \(K[G]\) 的中心）由元素 \(a = \sum a_g g\)（\(K[G]\) 中的元素）组成，这些元素满足 \(hah^{-1} = a\)（对每个 \(h \in G\)）。而今，我们有

\[
h a h ^ {- 1} = \sum_ {g \in \mathrm{G}} a _ {h ^ {- 1} g h} g;
\]

因此，当 G 有限时，中心 \(Z(K[G])\)（代数 K{[}G{]} 的中心）由中心函数组成。

设 \((\mathrm{M},\pi)\) 是 G 的一个线性表示。假设 M 是一个有限维自由 K-模。\(\pi\) 的迹是 K{[}G{]} 的表示（与 \(\pi\) 相伴的表示）的迹，即（VIII, 第 382 页）线性型 \(a\mapsto\mathrm{Tr}(\pi(a))\)（K{[}G{]} 上的线性型）。这个线性型由映射 \(g\mapsto\mathrm{Tr}(\pi(g))\)（从 G 到 K 的映射）确定，我们把该映射称为表示 \(\pi\) 的特征标，记作 \(\chi_{\pi}\)。表示的特征标是一个中心函数（II, §4, No.~3, 第 273 页，命题 3）。

设 \(M\) 与 \(M'\) 是有限维自由 \(K\)-模。设 \(\pi\) 与 \(\pi'\) 分别是 \(G\) 在 \(M\) 与 \(M'\) 中的线性表示。那么 \(M \oplus M'\) 是一个有限维自由 \(K\)-模，并且由 III, §9, No.~2, 第 543 页的命题 1，我们有

\[
\chi_ {\pi \oplus \pi^ {\prime}} = \chi_ {\pi} + \chi_ {\pi^ {\prime}}.
\]

更一般地，设

\[
0 \longrightarrow \mathrm{M} ^ {\prime} \longrightarrow \mathrm{M} \longrightarrow \mathrm{M} ^ {\prime \prime} \longrightarrow 0
\]

是 K{[}G{]}-模的一个正合列，并设 \(\pi\)、\(\pi'\)、\(\pi''\) 分别是与 M、\(M'\)、\(M''\) 相伴的 G 的线性表示。假设 K-模 \(M'\) 与 \(M''\) 都是自由且有限维的。那么 M 亦然（II, §1, No.~11, 第 218 页，命题 21），并且我们有

\[
\chi_ {\pi} = \chi_ {\pi^ {\prime}} + \chi_ {\pi^ {\prime \prime}}.\tag{1}
\]

命题 1. —— 假设 K 是一个交换域。设 \(\pi\) 与 \(\pi'\) 是 G 的有限维半单 K-线性表示。

\begin{enumerate}
\def\labelenumi{\alph{enumi})}
\item
  若 \(\pi(g)\) 与 \(\pi'(g)\) 的特征多项式对每个 \(g \in G\) 都相同，则表示 \(\pi\) 与 \(\pi'\) 同构。
\item
  再设 K 的特征为 0。若特征标 \(\chi_{\pi}\) 与 \(\chi_{\pi'}\) 相等，则表示 \(\pi\) 与 \(\pi'\) 同构。
\end{enumerate}

断言 a) 由 VIII, 第 386 页的推论 1 得出；断言 b) 由 VIII, 第 384 页的推论的 a) 得出。

例. —— 1) G 的单位表示或平凡表示是表示 \((\mathrm{K}, \varepsilon)\)，其中 \(\varepsilon(g) = \operatorname{Id}_{\mathrm{K}}\) 对每个 \(g \in G\) 成立。它的特征标是取常值 1 的常值函数。

\begin{enumerate}
\def\labelenumi{\arabic{enumi})}
\setcounter{enumi}{1}
\tightlist
\item
  G 的（左）正则表示是表示 \(\gamma\)（G 在 K{[}G{]} 中的表示），由 \(\gamma(g)(x)=gx\)（对 \(g\in G\) 与 \(x\in K[G]\)）定义。它对应于代数 K{[}G{]} 的左正则表示（VIII, 第 375 页）。设群 G 是有限的。此时正则表示是有限次数的。对 G 的每个不等于单位元的元素 g，g 在 K{[}G{]} 中的左乘关于典范基的矩阵是一个无不动点置换的矩阵。因此我们有
\end{enumerate}

\[
\chi_ {\boldsymbol {\gamma}} (g) = \left\{ \begin{array}{l l} | \mathrm{G} | & \text { if   g   is   the   identity   element }, \\ 0 & \text { otherwise }. \end{array} \right.\tag{2}
\]

我们类似地定义 G 的右正则表示。G 的双正则表示是表示 \(\rho\)（G × G 在 K{[}G{]} 中的表示），由 \(\rho(g, g')(x) = gxg'^{-1}\)（对 \((g, g') \in \mathrm{G} \times \mathrm{G}\) 与 \(x \in \mathrm{K}[\mathrm{G}]\)）定义。

\begin{enumerate}
\def\labelenumi{\arabic{enumi})}
\setcounter{enumi}{2}
\item
  给定 G 的一个线性表示 \((\mathrm{M},\pi)\)，逆步表示或对偶表示 \(\pi^{\vee}\)（\(\pi\) 的逆步表示或对偶表示）是 G 在与 M 对偶的 K-模中的表示，由关系式 \(\pi^{\vee}(g)={}^{t}\pi(g^{-1})\)（对每个 \(g\in G\)）定义（参见 II, §2, No.~5, 第 234 页）。若 M 是有限维自由 K-模，则它的对偶亦然，并且 \(\chi_{\pi^{\vee}}(g)=\chi_{\pi}(g^{-1})\) 对每个 \(g\in G\) 成立。
\item
  设 \((\mathrm{M},\pi)\) 与 \((\mathrm{M}^{\prime},\pi^{\prime})\) 是 G 的线性表示。在 VIII, 第 198 页的例 1 中，我们定义了一个 K{[}G{]}-模结构（在 \(M\otimes_{K}M^{\prime}\) 上）。相应的线性表示称为 \(\pi\) 与 \(\pi^{\prime}\) 的张量积，记作 \(\pi\otimes\pi^{\prime}\)。对 \(g\in G\)、\(x\in M\) 与 \(x^{\prime}\in M^{\prime}\)，有 \(\big((\pi\otimes\pi^{\prime})(g)\big)(x\otimes x^{\prime})=\pi(g)x\otimes\pi^{\prime}(g)x^{\prime}\)。
\end{enumerate}

若 M 与 M' 是有限维自由 K-模，则由 III, §9, No.~2, 第 543 页的命题 2，我们有

\[
\chi_ {\pi \otimes \pi^ {\prime}} = \chi_ {\pi} \chi_ {\pi^ {\prime}}.\tag{3}
\]

\begin{enumerate}
\def\labelenumi{\arabic{enumi})}
\setcounter{enumi}{4}
\tightlist
\item
  设 G 是两个群的乘积 \(G' \times G''\)。从 \(K[G'] \otimes_{K} K[G'']\) 到 K{[}G{]} 的、把 \(g' \otimes g''\) 映为 \((g', g'')\)（对 \(g' \in G'\) 与 \(g'' \in G''\)）的 K-线性映射是一个代数同构。设 \((M', \pi')\) 是 \(G'\) 的一个线性表示，\((M'', \pi'')\) 是 \(G''\) 的一个线性表示。\(\pi'\) 与 \(\pi''\) 的外张量积记作 \(\pi' \boxtimes \pi''\)，它是 \(G' \times G''\) 在向量空间 \(M' \otimes M''\) 中的表示，由 \((\pi' \boxtimes \pi'')(g', g'') = \pi'(g') \otimes \pi''(g'')\)（对 \((g', g'') \in G' \times G''\)）定义。若 \(M'\) 与 \(M''\) 是有限维自由 K-模，则 \(M' \otimes_{K} M''\) 是有限维自由 K-模，且表示 \(\pi' \boxtimes \pi''\) 的特征标由公式
\end{enumerate}

\[
\chi_ {\pi^ {\prime} \boxtimes \pi^ {\prime \prime}} (g ^ {\prime}, g ^ {\prime \prime}) = \chi_ {\pi^ {\prime}} (g ^ {\prime}) \chi_ {\pi^ {\prime \prime}} (g ^ {\prime \prime})
\]

（对 \(g' \in G'\) 与 \(g'' \in G''\)；III, §9, No.~2, 第 542 页，命题 2）。

\begin{enumerate}
\def\labelenumi{\arabic{enumi})}
\setcounter{enumi}{5}
\tightlist
\item
  设 \((V,\pi)\) 是 G 的一个线性表示，且 V 是有限维自由 K-模。我们用公式定义一个表示 \(\rho\)（\(G\times G\) 在 \(\operatorname{End}_{K}(V)\) 中的表示）：
\end{enumerate}

\[
\rho (g, g ^ {\prime}) (u) = \pi (g ^ {\prime}) \circ u \circ \pi (g ^ {- 1}).
\]

II, §4, No.~2, 第 271 页的 K-模同构 \(\theta_{\mathrm{V}}: \mathrm{V}^{*} \otimes_{\mathrm{K}} \mathrm{V} \to \mathrm{End}_{\mathrm{K}}(\mathrm{V})\) 是从外张量积 \(\pi^{\vee} \boxtimes \pi\) 到 \(\rho\) 的一个表示同构。映射 \(g \mapsto \rho(1, g)\)（相应地，\(g \mapsto \rho(g, 1)\)）是 G 的一个表示，它同构于 \(\pi^{\dim_{\mathrm{K}} \mathrm{V}}\)（相应地，\((\pi^{\vee})^{\dim_{\mathrm{K}} \mathrm{V}}\)）。

设 L 是一个交换 K-代数，\((\mathrm{M},\pi)\) 是群 G 的一个线性表示。群同态 \(\pi_{(\mathrm{L})}:\mathrm{G}\to\mathbf{GL}(\mathrm{M}_{(\mathrm{L})})\)（由 \(g\mapsto\mathrm{Id}_{\mathrm{L}}\otimes\pi(g)\) 定义）是 G 在 L-模 \(\mathrm{M}_{(\mathrm{L})}\) 中的一个线性表示，称为由表示 \(\pi\) 经标量环从 K 扩张到 L 而得到的 G 的线性表示。

设 K 是一个体，L 是一个非零交换 K-代数。设 \((M,\pi)\) 与 \((M',\pi')\) 是 G 的线性表示。表示 \(\pi\) 与 \(\pi'\) 同构，当且仅当 \(\pi_{(L)}\) 与 \(\pi'_{(L)}\) 同构（VIII, 第 37 页，定理 3）。

再设代数 L 是 K 的一个扩张。考虑 VIII, 第 198 页例 1 中定义的环 \(\mathrm{R}_{\mathrm{K}}(\mathrm{G})\) 与 \(\mathrm{R}_{\mathrm{L}}(\mathrm{G})\)。标量扩张定义了一个环同态

\[
u: \mathrm{R} _ {\mathrm{K}} (\mathrm{G}) \longrightarrow \mathrm{R} _ {\mathrm{L}} (\mathrm{G}).
\]

这个同态是单射，且元素 \(\xi \in \mathrm{R}_{\mathrm{K}}(\mathrm{G})\) 是有效的，当且仅当 \(u(\xi)\) 是有效的（VIII, 第 195 页，定理 1）。

